\chapter{Rheumatoid Factor (RF)}

\section*{What Is This Test?}
Rheumatoid factor (RF) is an autoantibody, most commonly of the IgM isotype, directed against the Fc portion of IgG. RF was the first autoantibody discovered in rheumatoid arthritis and has been used as a diagnostic and prognostic marker for RA since the 1940s. RF is produced by B cells in the synovium, lymph nodes, and bone marrow in response to immune complex formation and may play a role in the pathogenesis of RA by forming immune complexes that activate complement and recruit inflammatory cells \cite{smolen2016rheumatoid}. RF testing is performed using latex agglutination, nephelometry, or enzyme-linked immunosorbent assay (ELISA). It is important to note that RF is not specific for RA --- it is found in other autoimmune diseases (Sjogren syndrome, systemic lupus erythematosus), chronic infections (hepatitis C, subacute bacterial endocarditis, tuberculosis), and elderly individuals without disease.

\section*{Alternative Names}
RF; Rheumatoid Factor; Rheumatoid Arthritis Factor; RA Factor; RhF; Latex Fixation Test; RA Latex

\section*{Principle}
RF is measured using immunoturbidimetric or immunonephelometric assays on automated chemistry analyzers. In these methods, latex particles coated with aggregated human or rabbit IgG are incubated with patient serum. If RF is present, it binds to the Fc region of IgG, causing agglutination of the latex particles. The resulting increase in turbidity (light scattering) is measured photometrically and is proportional to the RF concentration. Older semi-quantitative methods include the latex agglutination slide test and the Waaler-Rose test (sheep red blood cell agglutination). ELISA methods can differentiate IgM, IgG, and IgA isotypes of RF, though IgM-RF is the most commonly measured isotype in clinical practice. Results are reported in IU/mL with calibration against the WHO International Reference Preparation for Rheumatoid Arthritis Serum.

\section*{History}
Rheumatoid factor was discovered in 1940 by Erik Waaler, a Norwegian physician, who observed that serum from patients with rheumatoid arthritis caused agglutination of sheep red blood cells sensitized with rabbit anti-sheep antibodies. This finding was independently confirmed by Harry Rose in 1948 (the ``Waaler-Rose test''). The latex fixation test using latex particles coated with human IgG was developed by Jacques Singer and Charles Plotz in 1956, greatly simplifying RF testing. The introduction of nephelometry and turbidimetry in the 1970s--1980s enabled quantitative automated measurement. The 1987 American College of Rheumatology classification criteria for RA included RF as one of seven criteria. The 2010 ACR/EULAR classification criteria retain RF (along with anti-CCP antibody) as a serologic marker, with high-positive RF (>3 times upper limit of normal) receiving greater weight \cite{aletaha2010}. In recent decades, the discovery of anti-citrullinated peptide antibodies (anti-CCP) has complemented RF testing, as anti-CCP offers superior specificity for RA (95--98\% vs. 75--80\% for RF).

\section*{Why This Test Is Ordered}
\begin{itemize}
  \item Diagnosis of rheumatoid arthritis (RA) in patients presenting with symmetric polyarthritis, morning stiffness lasting >1 hour, and inflammatory joint pain \cite{scott2010rheumatoid}
  \item Classification of RA according to the 2010 ACR/EULAR criteria, where RF positivity contributes to the serology score \cite{aletaha2010}
  \item Prognostic stratification of RA --- high-titer RF at diagnosis predicts more aggressive disease with erosive joint damage and extra-articular manifestations (rheumatoid nodules, vasculitis, interstitial lung disease)
  \item Differentiation of seropositive RA from seronegative RA (approximately 70--80\% of patients with established RA are seropositive)
  \item Evaluation of suspected Sjogren syndrome in patients with sicca symptoms and elevated RF
  \item Investigation of mixed cryoglobulinemia, particularly in hepatitis C virus (HCV) infection
  \item Evaluation of patients with palpable purpura, peripheral neuropathy, and renal dysfunction suspicious for cryoglobulinemic vasculitis
\end{itemize}

\section*{Primary vs Secondary Test}
This is a primary screening test for rheumatoid arthritis, used in conjunction with anti-CCP antibody testing. Positive RF supports the diagnosis of RA but must be interpreted in a complete clinical context as RF has limited specificity.

\section*{How This Test Is Performed}
For most patients, a health care professional will take a blood sample from a vein in your arm using a small needle. After the needle is inserted, a small amount of blood will be collected into a test tube or vial. You may feel a little sting when the needle goes in or out. This usually takes less than five minutes.

\section*{What Preparation Is Needed}
No special preparation is required. Fasting is not necessary. Patients should inform their healthcare provider about all medications, particularly anti-TNF agents and other biologic disease-modifying antirheumatic drugs (DMARDs), which may lower RF levels.

\section*{Contraindications and Precautions}
\begin{itemize}
  \item No absolute contraindications. Specimen should be collected in a serum separator tube (SST, gold-top) or red-top tube. Grossly hemolyzed or lipemic specimens may interfere with nephelometric measurements. RF positivity increases with age --- up to 15--25\% of healthy adults over age 65 have a positive RF without clinical disease \cite{ingegnoli2013rheumatoid}. False-positive RF can occur in chronic infections (subacute bacterial endocarditis, hepatitis C, tuberculosis, syphilis), chronic inflammatory conditions (sarcoidosis, interstitial lung disease), malignancies, and after vaccination. Cryoglobulin-containing samples may give falsely low results because cryoprecipitable RF may be lost during specimen processing; blood should be drawn and maintained at 37 degrees C if cryoglobulinemia is suspected.
\end{itemize}
\section*{Risks}
Minimal risk. Slight pain or bruising at the venipuncture site. Rare: hematoma, infection, vasovagal syncope.

\section*{What Happens After the Test}
Results are typically available within 1 to 2 hours. Positive RF results should be correlated with clinical findings and anti-CCP antibody testing. A positive RF alone (without anti-CCP) in the setting of symmetric arthritis should still prompt consideration of RA, but other causes of RF positivity must be excluded. Serial RF measurements are not recommended for monitoring disease activity in RA --- clinical assessment (joint counts, patient global assessment) and acute phase reactants (CRP, ESR) are preferred for this purpose.

\section*{Understanding Results}

\subsection*{Below Normal (Negative)}
\begin{itemize}
  \item Seronegative rheumatoid arthritis: Approximately 20--30\% of patients with RA are persistently RF-negative despite clinically and radiographically active disease. These patients may have a less aggressive disease course
  \item Early rheumatoid arthritis: RF may be negative in the first 6--12 months of disease before seroconversion occurs; anti-CCP antibody is more likely to be positive in early RA
  \item Seronegative spondyloarthropathies (ankylosing spondylitis, psoriatic arthritis, reactive arthritis, enteropathic arthritis) are classically RF-negative
  \item Osteoarthritis, gout, and calcium pyrophosphate deposition disease (pseudogout) are associated with negative RF
  \item RF positivity may be transient in the setting of acute infection and later become negative
\end{itemize}

\subsection*{Above Normal (Positive)}
\begin{itemize}
  \item Rheumatoid arthritis: RF is positive in approximately 70--80\% of patients with established RA. High-titer RF (>3 times the upper limit of normal) is associated with erosive joint disease, rheumatoid nodules, systemic vasculitis, and poorer long-term outcomes
  \item Sjogren syndrome: RF is positive in 60--90\% of patients with primary Sjogren syndrome. RF positivity in a patient with sicca symptoms, elevated anti-SSA/Ro, and anti-SSB/La supports the diagnosis
  \item Mixed cryoglobulinemia (type II and III): Typically associated with chronic hepatitis C virus infection. RF is positive in >90\% of patients. RF-positive cryoglobulinemia presents with palpable purpura, peripheral neuropathy, arthralgias, and membranoproliferative glomerulonephritis
  \item Systemic lupus erythematosus: RF is positive in 20--40\% of SLE patients and is more common in those with joint involvement
  \item Chronic infections: Subacute bacterial endocarditis (positive RF in 30--50\%), tuberculosis, syphilis, leprosy, and chronic viral hepatitis (particularly HCV) can all produce RF from polyclonal B-cell activation
  \item Interstitial lung disease: RF is positive in a subset of patients with idiopathic pulmonary fibrosis and other interstitial lung diseases
  \item Sarcoidosis: RF may be positive due to polyclonal B-cell activation from chronic granulomatous inflammation
  \item Healthy individuals: Low-titer positive RF is found in 5--10\% of healthy young adults and up to 15--25\% of healthy elderly individuals without clinical disease
\end{itemize}

\section*{Normal Results}
\begin{itemize}
  \item \textbf{RF (IgM):} <15--20 IU/mL (laboratory-dependent cutoff)
  \item \textbf{Age >65 years:} Low-titer positive RF found in 15--25\% of healthy elderly individuals
  \item \textbf{Weak positive:} 15--40 IU/mL
  \item \textbf{Moderate positive:} 40--60 IU/mL
  \item \textbf{Strong positive:} >60 IU/mL (>3 times upper limit of normal); this level contributes 3 points to the 2010 ACR/EULAR RA classification criteria serology score
  \item RF results should always be interpreted alongside anti-CCP antibody, clinical symptoms, and acute phase reactants (CRP, ESR)
\end{itemize}

\section*{Follow-up Tests}
\begin{itemize}
  \item \textbf{Anti-CCP (cyclic citrullinated peptide) antibody:} Higher specificity for RA (95--98\%) compared to RF. Anti-CCP can be positive in early RA and may precede clinical symptoms by years. Positive anti-CCP with positive RF confers the highest risk for erosive disease \cite{niewold2007anticcp}
  \item \textbf{Antinuclear antibody (ANA):} If features of systemic lupus erythematosus or other connective tissue disease are present
  \item \textbf{Anti-SSA/Ro and anti-SSB/La antibodies:} If Sjogren syndrome is suspected in a patient with sicca symptoms and positive RF
  \item \textbf{C-reactive protein (CRP) and erythrocyte sedimentation rate (ESR):} Acute phase reactants that reflect systemic inflammation in active RA. Used to monitor treatment response
  \item \textbf{Hepatitis C virus serology (anti-HCV):} If mixed cryoglobulinemia is suspected in a patient with positive RF, palpable purpura, and peripheral neuropathy
  \item \textbf{Radiographs of hands, wrists, and feet:} To assess for erosive changes, joint space narrowing, and periarticular osteopenia characteristic of RA
  \item \textbf{Complete blood count:} To evaluate for anemia of chronic disease commonly seen in active RA
  \item \textbf{Synovial fluid analysis:} If joint aspiration is performed to evaluate inflammatory versus non-inflammatory arthritis and to exclude crystal arthropathy or infection
\end{itemize}

\section*{References}
\bibliographystyle{unsrtnat}
\bibliography{references}

\clearpage
\section*{Regulatory Status and Authorization Date}
This chapter describes a test category, not a specific commercial product. FDA, CDSCO, EU IVDR, and MHRA approval or registration dates therefore cannot be assigned without the assay manufacturer, product name, instrument, and intended use. For laboratory-developed tests, an individual agency approval date may not exist. See the Preface for the interpretation of regulatory status.

\clearpage
